% ================================================================
\chapter{Character Creation}
% ================================================================
% REAL RULE: four columns, rated 1–4. Keep this.
\lettrine{E}{very character} is described by four columns, each rated on
a scale of \keyword{1 to 4} --- from ordinary to extraordinary: they are
\keyword{Body}, \keyword{Stuff}, \keyword{Skills}, and \keyword{Magic}.
Their meaning bends to fit each realm, but the rank you assign travels
with you from world to world.

\begin{statblock}[The Four Columns at a Glance]
\begin{tabularx}{\linewidth}{@{}p{0.18\linewidth} Y@{}}
\keyword{Body}   & Your physical form --- from ordinary person to dragon.\\[3pt]
\keyword{Stuff}  & What you carry --- from the shirt on your back to a ship.\\[3pt]
\keyword{Skills} & Everything mundane and learned, including combat.\\[3pt]
\keyword{Magic}  & Whatever makes the setting feel fantastic.\\
\end{tabularx}
\end{statblock}

\section{Spending Your Ranks}
% REAL RULE: Shadowrun-style priority. Assign one of A/B/C/D per column.
Character creation borrows the priority method from \emph{Shadowrun}.
You have four priorities --- \keyword{A}, \keyword{B}, \keyword{C}, and
\keyword{D} --- and you assign \emph{one letter to each column}, with no
repeats. \keyword{A} is where you shine; \keyword{D} is where you
struggle. No one is exceptional at everything.

The \keyword{Priority Matrix} on the following page lays out what each
letter means for each column. A priority of \keyword{A} grants the top
rank (\rating{4}) in that column, down to \keyword{D} for the lowest
(\rating{1}).

% ---- Full-width priority matrix: leads a fresh page, then columns resume.
% \twocolumn[<top material>] guarantees flush-top, full-width placement
% (a table* float would centre on a float-only page under memoir).
% NB: no bracketed \\[..] anywhere in this optional arg — a literal ]
% would terminate \twocolumn's [...] early. Use \par + \vspace instead.
\twocolumn[%
  \begingroup\centering
  {\hdfont\Large\bfseries\color{\realmaccent} The Priority Matrix\par}
  \vspace{2pt}
  {\itshape Assign one of A, B, C, D to each column --- one letter each, no repeats.\par}
  \vspace{10pt}
  \endgroup
  \begingroup
  \setlength{\tabcolsep}{7pt}\renewcommand{\arraystretch}{1.35}
  \noindent\begin{tabularx}{\textwidth}{@{}>{\centering\arraybackslash}p{2.3cm} *{4}{Y}@{}}
  \rowcolor{\realmaccent}
  {\hdfont\bfseries\color{white}Priority} &
  {\hdfont\bfseries\color{white}Body}   &
  {\hdfont\bfseries\color{white}Stuff}  &
  {\hdfont\bfseries\color{white}Skills} &
  {\hdfont\bfseries\color{white}Magic}  \\
  \rowcolor{atrBoxbg}
  {\hdfont\LARGE\bfseries\color{\realmaccent}A}\newline\rating{4}
    & Truly strange or monstrous --- a dragon, an angel.
    & A ship, a keep, incredible gear. Rich.
    & A master of many things.
    & A \emph{fucking wizard}; choose 12 tags. \\
  \rowcolor{white}
  {\hdfont\LARGE\bfseries\color{\realmaccent}B}\newline\rating{3}
    & Exceptional --- elven grace, ogrish might, uncanny beauty.
    & A legendary weapon or a vehicle; well off (\$20{,}000).
    & A master of your craft, or a jack of all trades.
    & You know magic; choose 7 tags. \\
  \rowcolor{atrBoxbg}
  {\hdfont\LARGE\bfseries\color{\realmaccent}C}\newline\rating{2}
    & You stand out --- strong, quick, or striking.
    & Ordinary things and some money (\$4{,}000).
    & You know a lot.
    & A few spells; choose 3 tags. \\
  \rowcolor{white}
  {\hdfont\LARGE\bfseries\color{\realmaccent}D}\newline\rating{1}
    & An ordinary person; base stats across the board.
    & Nothing but the shirt on your back.
    & You know a few things.
    & You do not know magic. \\
  \end{tabularx}
  \par\bigskip
  \endgroup
]

\begin{example}
% REAL worked example of the ABCD method.
Mara is a wandering swordmaster who has never touched real magic. She
assigns \keyword{A} to Skills, \keyword{B} to Body, \keyword{C} to Stuff,
and \keyword{D} to Magic.
\end{example}

\section{Stepping Through the Door}
% TODO: how play begins; hand-off to the column chapters.
\lipsum[9]

\begin{notebox}[Design Note: One Sheet, Many Worlds]
% Concept is real (rank travels, meaning is local); prose to be written.
Rank travels but meaning is local, so a single character sheet carries
you across every genre. \lipsum[10][1-3]
\end{notebox}
